\begin{textsample}{POS Dim 3 – es – Score 14.00 – es/es\_inf\_17.txt}  \label{ex:f3_pos_068}
Para Refletir Na primeira infância, o corpo é o alicerce para o desenvolvimento mental e emocional da \textbf{criança}, essencial na construção de \textbf{afetos} e sentimentos.
As experiências com o corpo, os \textbf{gestos} e os movimentos, promovidas para elas, constituem uma linguagem vital que orienta para o mundo.
Essas experiências devem ser ricas e plurais e serem promovidas de diversas formas : \textbf{gestos}, \textbf{mímicas}, posturas, movimentos expressivos levando as \textbf{crianças} a expressarem suas emoções, reconhecerem suas sensações, interagirem, \textbf{brincarem}, ocuparem os espaços localizando -se neles, construindo conhecimento de si e do mundo.
O corpo é a ferramenta ímpar usada pelas \textbf{crianças} para expressar -se e comunicar -se com excelência.
A capacidade de nomear, identificar e ter consciência do próprio corpo e a construção positiva da autoimagem favorecem a expressão e o conhecimento sobre a cultura corporal do meio em que vivem.
Quando \textbf{brincam}, dançam, \textbf{dramatizam}, as \textbf{crianças} expressam toda sua personalidade em construção, agem e imaginam situações, principalmente quando querem ter seus interesses atendidos.
Pensar este campo de experiências no trabalho pedagógico é fazer o exercício \textbf{diário} de garantir que as \textbf{crianças} tenham esses objetivos de aprendizagem e desenvolvimento efetivados.
Para tanto, cada docente deve pensar se : - Promove experiências diversas em que as \textbf{crianças} exploram os espaços com o corpo, potencializando suas habilidades? - Desenvolve jogos corporais levando -as a explorarem as formas básicas dos movimentos, suas dinâmicas e a ocupação desses espaços? - Permite as \textbf{brincadeiras} de jogos simbólicos evidenciadas por meio do faz de conta, momento em que as \textbf{crianças} representam os papéis sociais do seu cotidiano e o mundo da fantasia? - Planeja atividades com músicas, considerando a diversidade local, regional, nacional e internacional, permitindo as \textbf{crianças} recriarem \textbf{livremente} seus movimentos, ao mesmo tempo em que interagem com seus pares? - Uma vez que a linguagem cênica está integrada a muitas experimentações vivenciadas pelas \textbf{crianças} na Educação \textbf{Infantil}, como você tem organizado em sua \textbf{rotina} atividades como leitura literária, atividades com as artes visuais? - As \textbf{crianças} participam de apresentações teatrais?
O foco neste campo é que \textbf{bebês}, \textbf{crianças} bem \textbf{pequenas} e \textbf{crianças} \textbf{pequenas} explorem todas as possibilidades de se expressar, comunicar, interagir com seus parceiros, desenvolvendo \textbf{progressivamente} sua consciência corporal ( habilidades gestuais, posturas e movimentos ) com confiança, eficácia e autonomia.

% matched lemmas: afeto, bebê, brincadeira, brincar, criança, diário, dramatizar, gesto, infantil, livremente, mímico, pequeno, progressivamente, rotina
\end{textsample}
